%!TEX program = xelatex
\documentclass[a4paper,11pt]{article}

\usepackage{kotex}
\usepackage[margin=2.5cm, top=3cm, bottom=3cm]{geometry}
\usepackage{xcolor}
\usepackage[most]{tcolorbox}
\usepackage{enumitem}
\usepackage{booktabs}
\usepackage{titlesec}
\usepackage{fancyhdr}
\usepackage{hyperref}
\usepackage{fontspec}
\usepackage{multirow}

% ── 색상 정의 ──────────────────────────────────────────────────────
\definecolor{primary}{RGB}{37, 99, 235}      % blue-600
\definecolor{green}{RGB}{22, 163, 74}        % green-600
\definecolor{amber}{RGB}{217, 119, 6}        % amber-600
\definecolor{red}{RGB}{220, 38, 38}          % red-600
\definecolor{indigo}{RGB}{79, 70, 229}       % indigo-600
\definecolor{lightblue}{RGB}{239, 246, 255}  % blue-50
\definecolor{lightgreen}{RGB}{240, 253, 244} % green-50
\definecolor{lightamber}{RGB}{255, 251, 235} % amber-50
\definecolor{lightred}{RGB}{254, 242, 242}   % red-50
\definecolor{lightgray}{RGB}{249, 250, 251}  % gray-50

% ── 섹션 스타일 ────────────────────────────────────────────────────
\titleformat{\section}
  {\Large\bfseries\color{primary}}
  {\thesection.}{0.5em}{}[\vspace{-0.3em}\color{primary}\rule{\linewidth}{1pt}]

\titleformat{\subsection}
  {\normalsize\bfseries\color{black}}
  {\thesubsection}{0.5em}{}

% ── 박스 스타일 ────────────────────────────────────────────────────
\tcbset{
  tipbox/.style={
    colback=lightblue, colframe=primary, fonttitle=\bfseries\small,
    title={도움말}, arc=3pt, boxrule=0.5pt,
    left=6pt, right=6pt, top=5pt, bottom=5pt
  },
  warnbox/.style={
    colback=lightamber, colframe=amber, fonttitle=\bfseries\small,
    title={주의}, arc=3pt, boxrule=0.5pt,
    left=6pt, right=6pt, top=5pt, bottom=5pt
  },
  dangerbox/.style={
    colback=lightred, colframe=red, fonttitle=\bfseries\small,
    title={경고}, arc=3pt, boxrule=0.5pt,
    left=6pt, right=6pt, top=5pt, bottom=5pt
  },
  stepbox/.style={
    colback=lightgreen, colframe=green, arc=3pt, boxrule=0.5pt,
    left=6pt, right=6pt, top=5pt, bottom=5pt
  },
  infobox/.style={
    colback=lightgray, colframe=gray, arc=3pt, boxrule=0.5pt,
    left=6pt, right=6pt, top=5pt, bottom=5pt
  }
}

% ── 헤더/푸터 ──────────────────────────────────────────────────────
\pagestyle{fancy}
\fancyhf{}
\fancyhead[L]{\small\color{gray} 양현고 수학여행 앱}
\fancyhead[R]{\small\color{gray} 관리자 매뉴얼}
\fancyfoot[C]{\small\thepage}
\renewcommand{\headrulewidth}{0.4pt}

% ── 목록 설정 ──────────────────────────────────────────────────────
\setlist[enumerate]{leftmargin=1.5em, itemsep=2pt}
\setlist[itemize]{leftmargin=1.5em, itemsep=2pt, label={\textcolor{primary}{$\bullet$}}}

% ──────────────────────────────────────────────────────────────────
\begin{document}

% ── 표지 ──────────────────────────────────────────────────────────
\begin{titlepage}
  \centering
  \vspace*{3cm}
  {\Huge\bfseries\color{primary} 양현고 수학여행 앱}\par\vspace{0.5cm}
  {\LARGE 관리자 매뉴얼}\par\vspace{1cm}
  \textcolor{gray}{\rule{8cm}{0.5pt}}\par\vspace{1cm}
  {\large 2026 수학여행 현장 관리 시스템}\par\vspace{0.5cm}
  {\normalsize\color{gray} v1.0 · 2026년 5월}\par
  \vfill
  \begin{tcolorbox}[colback=lightblue, colframe=primary, arc=6pt,
    boxrule=1pt, width=10cm]
    \centering
    \textbf{접속 주소}\par\vspace{0.3em}
    {\large\color{primary}\texttt{yanghyeon-tour-2026.vercel.app}}
  \end{tcolorbox}
\end{titlepage}

\tableofcontents
\newpage

% ──────────────────────────────────────────────────────────────────
\section{관리자 계정과 접속}

\subsection{관리자 계정}

관리자 계정은 학교 서버 측에서 직접 발급합니다. Firebase Console에서 계정을 만들고
Custom Claim \texttt{role: "admin"}을 부여해야 합니다. 일반 가입 절차로는 관리자 권한을
취득할 수 없습니다.

\subsection{로그인 및 화면 진입}

\begin{enumerate}
  \item 브라우저에서 \textbf{\texttt{yanghyeon-tour-2026.vercel.app}} 접속
  \item 관리자 이메일과 비밀번호로 로그인
  \item 자동으로 \textbf{/admin} 경로(관리자 페이지)로 이동
\end{enumerate}

\begin{tcolorbox}[warnbox]
교사 또는 학생 계정으로 로그인하면 각 역할의 포털로 자동 이동됩니다.
관리자 페이지(\texttt{/admin})에 직접 접근을 시도해도 권한 확인 후 리디렉션됩니다.
\end{tcolorbox}

\subsection{화면 구조}

관리자 페이지는 상단 헤더와 5개의 탭으로 구성됩니다.

\begin{center}
\begin{tabular}{@{}lll@{}}
  \toprule
  \textbf{탭} & \textbf{기능 요약} & \textbf{배지} \\
  \midrule
  대시보드 & 통계, 운영 액션, 차트 & \\
  학생 & 전체 학생 목록, 가입 상태 조회 & 학생 수 \\
  교사 & 가입된 교사 계정 목록 & \\
  점호 & 세션 생성·연장·종료 & 활성 세션 수 \\
  설정 & 전역 설정, UID 초기화 & \\
  \bottomrule
\end{tabular}
\end{center}

% ──────────────────────────────────────────────────────────────────
\section{대시보드 탭}

대시보드는 관리자가 가장 자주 사용하는 탭으로, 현황 파악과 핵심 운영 기능을 모두 담고 있습니다.

\subsection{진행 중 점호 배너}

활성 점호 세션이 있을 때 대시보드 최상단에 \textbf{주황색 배너}가 나타납니다.

\begin{itemize}
  \item 세션 이름, 유형(정시점호/승차점호), 범위 표시
  \item 복수 세션이 있으면 \texttt{+N건 더}로 추가 건수 표시
  \item 배너를 탭하면 해당 점호 모니터링 화면(\texttt{/teacher/checkin})으로 이동
\end{itemize}

\subsection{통계 카드 (4종)}

\begin{center}
\begin{tabular}{@{}lll@{}}
  \toprule
  \textbf{카드} & \textbf{표시 내용} & \textbf{색상} \\
  \midrule
  전체 학생 & 전체 인원 수 / 앱 가입자 수 & 파랑 \\
  진행 점호 & 현재 활성 세션 수 & 주황 \\
  미처리 사건 & 종결되지 않은 사건사고 수 & 빨강(건수 있을 때) \\
  누적 사건 & 전체 등록 사건사고 수 & 회색 \\
  \bottomrule
\end{tabular}
\end{center}

\subsection{운영 액션}

\subsubsection{공지 일괄 발송}

주황색 카드를 탭하면 다이얼로그가 열립니다.

\begin{tcolorbox}[stepbox]
\begin{enumerate}
  \item \textbf{대상} 선택: 전체 반 / 특정 학년 / 반 직접 선택
  \item 「특정 학년」 선택 시 → 학년 버튼 추가 선택
  \item 「반 직접 선택」 선택 시 → 학년별 반 버튼을 개별 선택
  \item \textbf{메시지} 입력 (최대 500자)
  \item 필요 시 「교직원 방에도 발송 사본 남기기」 체크
  \item \textbf{N개 반에 발송} 버튼 클릭
\end{enumerate}
\end{tcolorbox}

\begin{tcolorbox}[tipbox]
발송된 공지는 각 반의 공지방(채팅)에 메시지로 추가됩니다.
학생들은 앱에서 읽거나 푸시 알림으로 수신합니다.
\end{tcolorbox}

\subsubsection{사건사고 관리}

\begin{itemize}
  \item 미처리 사건 수가 표시된 빨간 카드 — 탭하면 \texttt{/teacher/incident}로 이동
  \item 해당 페이지에서 사건 등록, 상세 확인, 종결 처리 가능
\end{itemize}

\subsubsection{승차 현황 모니터링}

\begin{itemize}
  \item 파란 버스 카드 — 탭하면 \texttt{/teacher/boarding}으로 이동
  \item 전체 반 승차 현황을 버스 좌석 배치도로 확인
\end{itemize}

\subsection{현황 차트}

\begin{itemize}
  \item \textbf{가입·학년별 통계}: 학년별 가입/미가입 비율 막대 그래프
  \item \textbf{반별 학생 분포}: 전체 반 학생 수 분포 차트
\end{itemize}

\subsection{조회·소통 (바로가기 그리드)}

대시보드 하단의 4칸 그리드에서 바로 이동할 수 있습니다.

\begin{center}
\begin{tabular}{@{}lll@{}}
  \toprule
  \textbf{버튼} & \textbf{이동 경로} & \textbf{비고} \\
  \midrule
  채팅 & \texttt{/chat} & 공지방 열람 \\
  일정 & \texttt{/schedule} & 여행 일정 조회 \\
  연락처 & \texttt{/contacts} & 비상 연락처 \\
  점호 뷰 & \texttt{/teacher/checkin?session=latest} & 최신 세션 모니터링 \\
  \bottomrule
\end{tabular}
\end{center}

\subsection{데이터 관리}

\begin{center}
\begin{tabular}{@{}lll@{}}
  \toprule
  \textbf{버튼} & \textbf{기능} & \textbf{이동/동작} \\
  \midrule
  Excel 업로드 & 학생·호실·일정·버스·연락처 일괄 등록 & \texttt{/admin/upload} 이동 \\
  버스 QR 인쇄 & 호차별 QR 코드 PDF 생성 후 인쇄 & 브라우저 인쇄 다이얼로그 \\
  \bottomrule
\end{tabular}
\end{center}

\begin{tcolorbox}[tipbox]
\textbf{버스 QR 인쇄}: 버스 데이터가 업로드된 후에 사용하세요.
각 호차 QR 코드를 한 장씩 출력해 해당 버스에 부착합니다.
\end{tcolorbox}

% ──────────────────────────────────────────────────────────────────
\section{학생 탭}

전체 학생 명단을 조회하고 앱 가입 여부를 확인합니다.

\subsection{검색 및 필터}

\begin{itemize}
  \item \textbf{검색창}: 이름, 학년, 반으로 검색 가능
  \item \textbf{필터 버튼}: 전체 / 가입 / 미가입 전환
  \item 검색 결과 수가 오른쪽에 표시
\end{itemize}

\subsection{학생 카드 정보}

각 학생 카드에 다음 정보가 표시됩니다:

\begin{center}
\begin{tabular}{@{}ll@{}}
  \toprule
  \textbf{항목} & \textbf{내용} \\
  \midrule
  이름 & 학생 이름 \\
  학년·반·번호 & 소속 정보 \\
  호실·호차 & 숙소 방 번호 / 탑승 버스 번호 \\
  가입 배지 & 초록 「가입」 / 회색 「미가입」 \\
  \bottomrule
\end{tabular}
\end{center}

\begin{tcolorbox}[tipbox]
미가입 학생이 있는 경우 해당 학생에게 담임 교사를 통해 계정 정보를 전달하세요.
학생 가입 코드 발급은 설정 탭의 \textbf{교사 가입 코드} 항목을 참고하세요.
(학생 가입은 학번 형식 이메일 + 초기 비밀번호로 이루어지며, 별도 코드는 필요하지 않습니다.)
\end{tcolorbox}

% ──────────────────────────────────────────────────────────────────
\section{교사 탭}

앱에 가입된 교사 계정 목록을 확인합니다.

\begin{itemize}
  \item 가입 교사 수가 상단에 표시
  \item 각 카드에 \textbf{이름}과 \textbf{담임 학년·반} 표시
  \item 담임이 지정되지 않은 교사는 「담임 미지정」으로 표시
\end{itemize}

\begin{tcolorbox}[tipbox]
담임 정보 자동 설정은 Excel 업로드 페이지의 \textbf{담임 정보 동기화} 기능을 사용하세요.
버스 시트의 인솔교사1 이름을 기준으로 교사 계정에 담임반이 자동 매핑됩니다.
\end{tcolorbox}

% ──────────────────────────────────────────────────────────────────
\section{점호 탭}

점호 세션을 생성하고 현재 진행 중인 세션을 관리합니다.

\subsection{세션 생성}

오른쪽 상단 \textbf{「새 세션」} 버튼을 탭합니다.

\begin{tcolorbox}[stepbox]
\begin{enumerate}
  \item \textbf{점호 유형} 선택: 정시점호 (야간 취침) / 승차점호 (버스 탑승)
  \item \textbf{세션 이름} 입력 (선택): 예) 2일차 저녁 점호
  \item \textbf{범위} 선택
    \begin{itemize}
      \item 전체: 모든 학생
      \item 학급: 특정 반 번호 입력 (예: 1)
      \item 호실: 특정 방 번호 입력 (예: 201)
      \item 호차: 특정 버스 번호 입력 (예: 1)
    \end{itemize}
  \item \textbf{운영 시간} 선택: 10분 / 20분 / 30분 / 60분
  \item \textbf{「세션 시작」} 버튼 클릭
\end{enumerate}
\end{tcolorbox}

\begin{tcolorbox}[tipbox]
이름을 입력하지 않으면 「점호유형 HH:MM」 형식으로 자동 생성됩니다.\\
자동 점호(설정 탭 참고)를 사용하면 일정표 시각에 맞춰 세션이 자동 생성됩니다.
\end{tcolorbox}

\subsection{세션 관리}

진행 중인 각 세션 카드에서 할 수 있는 작업:

\begin{center}
\begin{tabular}{@{}lll@{}}
  \toprule
  \textbf{버튼} & \textbf{기능} & \textbf{비고} \\
  \midrule
  +10분 & 세션 종료 시각을 10분 연장 & 여러 번 클릭 가능 \\
  종료 & 세션 즉시 종료 & 종료 후 복구 불가 \\
  \bottomrule
\end{tabular}
\end{center}

\begin{itemize}
  \item 세션 카드 우상단 \textbf{초록 점}: 활성 (남은 시간 표시)
  \item \textbf{빨간 점}: 만료 (자동 종료됨)
\end{itemize}

\begin{tcolorbox}[warnbox]
교사는 세션 연장·종료 권한이 없습니다. (Firestore 규칙에 의해 DB 레벨에서 차단)\\
세션 관리는 반드시 관리자 계정으로 수행하세요.
\end{tcolorbox}

% ──────────────────────────────────────────────────────────────────
\section{설정 탭}

\subsection{전역 설정}

\subsubsection{교사 가입 코드}

교사가 앱에 가입할 때 입력해야 하는 인증 코드입니다.

\begin{itemize}
  \item 코드 입력 후 \textbf{「저장」} 버튼으로 즉시 적용
  \item 수학여행 전 담임 교사들에게 코드를 별도 전달하세요
  \item 수학여행 종료 후 코드를 변경하면 교사 계정 추가 가입을 차단할 수 있습니다
\end{itemize}

\subsubsection{학생 가입 잠금}

\begin{center}
\begin{tabular}{@{}lll@{}}
  \toprule
  \textbf{상태} & \textbf{아이콘} & \textbf{효과} \\
  \midrule
  잠금 해제 & 회색 토글 & 학생이 앱에 새로 가입 가능 \\
  잠금 & 주황색 토글 & 신규 학생 가입 불가 (기존 계정 로그인은 가능) \\
  \bottomrule
\end{tabular}
\end{center}

\begin{tcolorbox}[tipbox]
수학여행 출발 당일에는 잠금을 해제하여 미가입 학생이 현장에서 가입할 수 있도록 하세요.
출발 후에는 잠금하는 것을 권장합니다.
\end{tcolorbox}

\subsubsection{QR 승차점호 활성}

\begin{center}
\begin{tabular}{@{}lll@{}}
  \toprule
  \textbf{상태} & \textbf{아이콘} & \textbf{효과} \\
  \midrule
  활성 & 초록 토글 & 학생이 버스 QR 스캔으로 승차 확인 가능 \\
  비활성 & 회색 토글 & QR 스캔 차단 (수동 점호만 가능) \\
  \bottomrule
\end{tabular}
\end{center}

\subsubsection{자동 점호}

일정표에 등록된 시각이 되면 자동으로 점호 세션을 생성하는 기능입니다.

\begin{center}
\begin{tabular}{@{}lll@{}}
  \toprule
  \textbf{상태} & \textbf{아이콘} & \textbf{효과} \\
  \midrule
  활성 & 파랑 토글 & 일정 시각에 자동 세션 생성 \\
  비활성 & 회색 토글 & 자동 생성 없음 (수동 생성만) \\
  \bottomrule
\end{tabular}
\end{center}

\begin{tcolorbox}[warnbox]
\textbf{자동 점호는 관리자 페이지(브라우저 탭)가 열려 있을 때만 동작합니다.}\\
Cloud Functions를 사용하지 않으므로, 관리자 기기의 브라우저를 수학여행 기간 내내 켜 두어야 합니다.
탭을 닫거나 절전 상태가 되면 자동 점호가 실행되지 않습니다.
\end{tcolorbox}

\subsection{학생 UID 초기화}

학생의 앱 가입 정보를 초기화합니다. 초기화 후 해당 학생은 다시 가입할 수 있습니다.

\begin{tcolorbox}[stepbox]
\begin{enumerate}
  \item 검색창에 이름 또는 반을 입력하여 해당 학생 검색
  \item 대상 학생 오른쪽 \textbf{「초기화」} 버튼 클릭
  \item 확인 다이얼로그에서 내용 확인 후 \textbf{확인}
\end{enumerate}
\end{tcolorbox}

\begin{tcolorbox}[dangerbox]
\textbf{UID 초기화는 되돌릴 수 없습니다.}\\
초기화 시 Firestore의 학생 문서에서 uid 필드만 제거됩니다.
Firebase Authentication에 생성된 실제 계정은 별도로 Firebase Console에서 수동 삭제해야 합니다.\\
계정을 삭제하지 않으면 같은 학번으로 재가입을 시도할 때 오류가 발생할 수 있습니다.
\end{tcolorbox}

% ──────────────────────────────────────────────────────────────────
\section{Excel 업로드 페이지}

대시보드 「Excel 업로드」 카드를 탭하면 \texttt{/admin/upload}로 이동합니다.
수학여행 전 모든 데이터를 한 번에 Firestore에 등록하는 데 사용합니다.

\subsection{Excel 파일 형식}

\begin{tcolorbox}[warnbox]
Excel 파일은 반드시 \textbf{5개 시트}를 포함해야 합니다.
각 시트의 \textbf{1행이 컬럼 헤더}여야 합니다. (설명 행 없이 바로 헤더 시작)
\end{tcolorbox}

\begin{center}
\begin{tabular}{@{}lll@{}}
  \toprule
  \textbf{시트명} & \textbf{내용} & \textbf{비고} \\
  \midrule
  \texttt{students} & 학생 명단 & uid/createdAt 유지하며 upsert \\
  \texttt{rooms} & 호실 정보 & upsert \\
  \texttt{schedule} & 여행 일정 & 기존 데이터 삭제 후 재업로드 \\
  \texttt{buses} & 버스 정보 & upsert \\
  \texttt{contacts} & 비상 연락처 & 기존 데이터 삭제 후 재업로드 \\
  \bottomrule
\end{tabular}
\end{center}

\subsection{업로드 절차}

\begin{tcolorbox}[stepbox]
\begin{enumerate}
  \item 업로드 페이지에서 \textbf{「파일 선택」} 클릭
  \item \texttt{.xlsx} 또는 \texttt{.xls} 파일 선택
  \item \textbf{파싱 결과 미리보기}에서 각 시트별 인식된 건수 확인
  \item 일정표 진단 정보를 펼쳐 컬럼 인식 여부 확인 (필요 시 수정)
  \item \textbf{「Firestore에 업로드」} 클릭
  \item 시트별 업로드 진행 상태 확인 후 완료 메시지 확인
\end{enumerate}
\end{tcolorbox}

\subsection{일정표 컬럼 인식 규칙}

일정표 시트는 다음 컬럼명을 자동으로 인식합니다.

\begin{center}
\begin{tabular}{@{}ll@{}}
  \toprule
  \textbf{필드} & \textbf{인식 가능한 컬럼명} \\
  \midrule
  일차 & 일차, Day, 날짜, 일, 차수, 일자, 구분, day \\
  시작시각 & 시작시각, 시작, 시작시간, Start, 시작 시각, 시작 시간, start \\
  \bottomrule
\end{tabular}
\end{center}

인식 실패 시 해당 컬럼명을 위 목록 중 하나로 변경한 뒤 재업로드하세요.

\subsection{점호 유형 일괄 설정}

업로드 페이지 하단의 별도 카드입니다. 일정표 항목에 점호 유형을 일괄 적용합니다.

\begin{itemize}
  \item \textbf{정시점호} 적용 대상: 「학교 집결」, 「숙소 도착 후 휴식」 일정
  \item \textbf{승차점호} 적용 대상: 그 외 모든 일정
  \item \textbf{「Firestore 일괄 적용」} 버튼 클릭 시 즉시 반영
\end{itemize}

\begin{tcolorbox}[tipbox]
일정표 업로드 직후에 이 기능을 실행하면 점호 유형이 자동으로 분류됩니다.
분류 후 필요한 경우 Firebase Console에서 개별 수정도 가능합니다.
\end{tcolorbox}

\subsection{담임 정보 동기화}

교사 계정에 담임 학년·반을 자동으로 설정합니다.

\begin{itemize}
  \item 버스 시트의 \textbf{인솔교사1} 이름과 \textbf{탑승반}을 기준으로 매핑
  \item 이름이 정확히 일치하는 교사 계정에만 담임반이 설정됨
  \item \textbf{「담임 정보 동기화」} 버튼 클릭 시 업데이트된 교사 수 알림 표시
\end{itemize}

\begin{tcolorbox}[warnbox]
교사 이름이 버스 시트와 앱 가입 이름이 정확히 일치해야 합니다.
동기화 후 교사 탭에서 담임 지정 여부를 반드시 확인하세요.
\end{tcolorbox}

% ──────────────────────────────────────────────────────────────────
\section{수학여행 전 준비 체크리스트}

\begin{tcolorbox}[infobox]
\begin{enumerate}
  \item[\checkmark] Excel 파일 5개 시트 데이터 준비 완료
  \item[\checkmark] \texttt{/admin/upload}에서 Excel 업로드 완료
  \item[\checkmark] 점호 유형 일괄 설정 완료
  \item[\checkmark] 담임 정보 동기화 완료 (교사 탭에서 확인)
  \item[\checkmark] 교사 가입 코드 설정 후 담임 교사들에게 전달
  \item[\checkmark] 학생 가입 잠금 \textbf{해제} 상태 확인 (출발 당일까지)
  \item[\checkmark] QR 승차점호 활성 상태 확인
  \item[\checkmark] 버스 QR 인쇄 후 각 호차에 부착
  \item[\checkmark] 자동 점호 활성화 (관리자 기기 브라우저 상시 유지)
  \item[\checkmark] 출발 후 학생 가입 잠금 \textbf{활성화}
\end{enumerate}
\end{tcolorbox}

% ──────────────────────────────────────────────────────────────────
\section{자주 묻는 질문}

\begin{description}[leftmargin=0pt, style=nextline]
  \item[\textbf{Q. 교사가 점호 세션을 종료하려고 하는데 버튼이 없다고 합니다.}]
    정상적인 동작입니다. 점호 세션의 연장·종료는 관리자 전용 기능으로,
    교사 화면에는 해당 버튼이 표시되지 않습니다.
    관리자 계정의 점호 탭에서 처리하세요.

  \item[\textbf{Q. 자동 점호가 실행되지 않습니다.}]
    자동 점호는 관리자 페이지가 브라우저에 열려 있을 때만 동작합니다.
    설정 탭에서 자동 점호 토글이 활성(파랑)인지 확인하고,
    관리자 기기의 브라우저 탭이 열려 있는지 확인하세요.

  \item[\textbf{Q. Excel 업로드 후 일정표 점호 배지가 표시되지 않습니다.}]
    업로드 페이지에서 「일정 점호 유형 일괄 설정」의 \textbf{「Firestore 일괄 적용」}을
    실행했는지 확인하세요. 이 단계를 건너뛰면 점호 유형이 빈 값으로 남습니다.

  \item[\textbf{Q. 교사 탭에 담임이 미지정으로 표시됩니다.}]
    업로드 페이지의 「담임 정보 동기화」를 실행하세요.
    버스 시트의 인솔교사1 이름과 교사 앱 가입 이름이 동일해야 매핑됩니다.

  \item[\textbf{Q. 학생이 UID를 초기화했는데 재가입이 안 된다고 합니다.}]
    Firebase Console → Authentication에서 해당 학생의 계정을 삭제해야 합니다.
    UID 초기화는 Firestore의 연결만 끊을 뿐, Firebase Auth 계정은 삭제하지 않습니다.

  \item[\textbf{Q. QR 인쇄 버튼을 눌렀는데 PDF가 나오지 않습니다.}]
    버스 데이터가 Firestore에 업로드되어 있는지 확인하세요.
    버스 데이터가 없으면 「버스 데이터가 없습니다」 오류가 표시됩니다.

  \item[\textbf{Q. 공지 발송 후 학생이 알림을 받지 못했다고 합니다.}]
    학생 기기의 브라우저 알림 권한이 허용되어 있는지 확인하도록 안내하세요.
    알림이 오지 않더라도 앱의 공지방을 열면 메시지를 확인할 수 있습니다.
\end{description}

\end{document}
